\subsection{Robustness of the Relevant-Systematics List Under a Low-$E_{\mathrm{avail}}$ Cut}
\label{sec:relevantSystsQ0Lt100MeV}

Section~\ref{sec:relevantSystsDataDriven} derived the 26-systematic
neutron-multiplicity whitelist from diagnostics evaluated over the fully
inclusive sample. To check whether that list still holds in the
low-hadronic-energy ($E_{\mathrm{avail}}<100$~MeV) region used elsewhere in
this analysis, the same procedure -- $\pm1\sigma$ shifts compared to nominal
via $\chi^2_{\mathrm{shape}}=\sum_i(s_i-n_i)^2/n_i$ on the rescaled shape,
same $\chi^2_{\mathrm{shape}}>5$ threshold, same 77 systematics and 14
diagnostics -- was repeated with every diagnostic additionally restricted to
$E_{\mathrm{avail}}<100$~MeV. Only \textbf{11 of the 77} systematics clear
the threshold in this restricted sample, compared to 26 inclusively, and
every one of the 11 was already part of the inclusive whitelist -- the cut
never promotes a new systematic to relevant, it only suppresses
significance.

\begin{table}[htbp]
    \centering
    \caption{Systematics retained (of those evaluated per category), fully
    inclusive vs.\ restricted to $E_{\mathrm{avail}}<100$~MeV.}
    \label{table:relevantSystsCategoryCompare}
    \begin{tabular}{lccc}
        \toprule
        \textbf{Category} & \textbf{Inclusive} & \textbf{$E_{\mathrm{avail}}<100$~MeV} & \textbf{Change} \\
        \midrule
        FSI          & 5/5  & 0/5  & all dropped \\
        MEC          & 6/6  & 4/6  & $-2$ \\
        RES          & 6/10 & 3/10 & $-3$ \\
        QE           & 3/7  & 2/7  & $-1$ \\
        DIS          & 6/40 & 2/40 & $-4$ \\
        LowPriority  & 0/9  & 0/9  & unchanged \\
        \midrule
        \textbf{Total} & \textbf{26/77} & \textbf{11/77} & $-15$ \\
        \bottomrule
    \end{tabular}
\end{table}

The most striking change is FSI: every FSI knob drops below threshold,
led by \texttt{hNFSI\_MFP\_2024} (mean free path) falling from
$\chi^2_{\mathrm{shape}}=179.5$ to $2.0$ and \texttt{FormZone2024} from
$77.3$ to $0.2$. This is physically expected rather than a fluctuation: FSI
rescattering/absorption acts mainly by adding secondary-scattering hadronic
activity, i.e.\ by pushing events to \emph{higher} $E_{\mathrm{avail}}$, so
restricting to $E_{\mathrm{avail}}<100$~MeV removes almost exactly the phase
space where FSI has any effect. MEC, RES, QE, and DIS each keep roughly
half their inclusive members; a few (\texttt{MECShape2024AntiNu},
\texttt{LowQ2RESSupp2020}) actually grow \emph{more} significant under the
cut, since their shape effect is concentrated at low $E_{\mathrm{avail}}$
rather than diluted by it.

\begin{table}[htbp]
    \centering
    \caption{The 11 systematics retained under $E_{\mathrm{avail}}<100$~MeV,
    with $\chi^2_{\mathrm{shape}}$ in both samples for comparison. Sorted by
    the restricted-sample value. All 11 were already in the inclusive
    26-systematic list (Table~\ref{table:relevantSysts}).}
    \label{table:relevantSystsQ0Lt100MeV}
    \begin{tabular}{p{0.26\linewidth} p{0.08\linewidth} p{0.11\linewidth} p{0.11\linewidth} p{0.28\linewidth}}
        \toprule
        \textbf{Systematic} & \textbf{Cat.} & \textbf{$\chi^2$ incl.} & \textbf{$\chi^2$ $<$100 MeV} & \textbf{Diagnostic ($<$100 MeV, beam)} \\
        \midrule
        \texttt{MECShape2024AntiNu}           & MEC & 19.5  & 34.6 & final-state topology (RHC) \\
        \texttt{MECShape2024Nu}               & MEC & 66.4  & 30.1 & final-state topology (FHC) \\
        \texttt{RESDeltaScaleSyst}             & RES & 55.2  & 19.8 & final-state topology (FHC) \\
        \texttt{DISvnCC1pi\_2020}              & DIS & 166.5 & 19.4 & final-state topology (FHC) \\
        \texttt{DISNuHadronQ1Syst}             & DIS & 184.2 & 18.7 & final-state topology (FHC) \\
        \texttt{MECInitStateNPFrac2020Nu}      & MEC & 78.1  & 15.8 & final-state topology (FHC) \\
        \texttt{ZNormCCQE}                     & QE  & 104.2 & 13.3 & inclusive $E_{\mathrm{avail}}$ (RHC) \\
        \texttt{MECInitStateNPFrac2020AntiNu}  & MEC & 23.5  & 12.5 & final-state topology (RHC) \\
        \texttt{RPAShapesupp2020}              & QE  & 29.5  & 11.6 & inclusive $E_{\mathrm{avail}}$ (FHC) \\
        \texttt{LowQ2RESSupp2020}              & RES & 5.8   & 8.7  & final-state topology (FHC) \\
        \texttt{MaCCRES}                       & RES & 58.5  & 7.1  & final-state topology (FHC) \\
        \bottomrule
    \end{tabular}
\end{table}

For several of these (\texttt{MECShape2024Nu}, \texttt{MaCCRES},
\texttt{RESDeltaScaleSyst}, \texttt{LowQ2RESSupp2020}), the most sensitive
diagnostic also shifts from $E_{\mathrm{avail}}$ itself, inclusively, to
final-state topology once $E_{\mathrm{avail}}<100$~MeV is imposed --
unsurprising, since compressing $E_{\mathrm{avail}}$ into a narrow window
weakens its own shape-discriminating power, while the topology diagnostic
is unaffected by the cut. Overall, the neutron-multiplicity GENIE
multiverse built for the inclusive sample (Sec.~\ref{sec:relevantSystsDataDriven})
remains a conservative superset for a low-$E_{\mathrm{avail}}$-restricted
analysis: a multiverse built from just these 11 systematics is smaller and
cheaper, but every systematic it needs is already covered by the
inclusive 26.
